\documentclass[main.tex]{subfiles}


\begin{document}

%from pagenumber to up
% \phantomsection 
% \hypertarget{toc}{}

 

 

 
   

\section{\emph{p}--\emph{n}-Переход}


$p$--$n$-\emph{Переход}~--- это место соприкосновения $p$- и $n$-полупроводников (рис.\,\ref{fig:p132}).

% На рис.\,\ref{fig:p132} изображен контакт $p$- и $n$-полупроводников.

\tikzfading[name=fade right,
left color=transparent!0,
right color=transparent!100]

\tikzfading[name=fade left,
left color=transparent!100,
right color=transparent!0]

    \begin{figure}[H]
    \centering

    \begin{tikzpicture}[font=\small,            line cap=round,                             >={Stealth[inset=0pt,round,length=3mm, angle'=20]},vec/.style={>={Stealth[inset=0pt,round,length=3.5mm, angle'=20]}}]


\tikzset{el/.pic={
  \begin{scope}[shift={({rnd*360}:.05)}]
  \shade[ball color=NavyBlue!70] (0,0) circle (\de);
  \end{scope}}
}


\tikzset{pp/.pic={
  \begin{scope}[shift={({rnd*360}:.05)}]
  \draw[red,thick] (0,0) circle (\de);
  \end{scope}}
}


\tikzset{em/.pic={
  \begin{scope}[shift={({rnd*360}:.05)}]
  \begin{scope}[transparency group,nearly transparent]
  \begin{scope}[transparency group,semitransparent]
  \shade[ball color=NavyBlue!70] ([shift={(180:\de cm)}]0,0) circle (\de);
  \end{scope}
  \begin{scope}[transparency group]
  \shade[ball color=NavyBlue!70] ([shift={(180:{0 cm})}]0,0) circle (\de);
  \end{scope}
  \end{scope}
  \shade[ball color=NavyBlue!70] ([shift={(0:{\de cm})}]0,0) circle (\de);
  \end{scope}}
}


\def\dn{.2}
\def\de{.1}


\def\sr{3.5}
%p n
\fill[orange, nearly transparent] (0,0) rectangle++ (\sr,2);
\fill[Green, nearly transparent] (\sr,0) rectangle++ (\sr,2);


% border charges
\fill[red,path fading=fade right] (\sr,0) rectangle++ (.5,2);
\fill[NavyBlue,path fading=fade left] (\sr,0) rectangle++ (-.5,2);


% border
\draw[gray] (\sr,0) --++ (0,2);
\draw (0,0) rectangle++ (2*\sr,2);


%%%%%%%%%%%%%%%%%%%
% p solid 
\foreach \y in {0}
\foreach \x in {0,...,5}
\pic[shift={(.3,.3)}] at ({0+(3-.6)/5*\x},{0+(2-.6)/3*\y}) {pp};

\foreach \y in {1}
\foreach \x in {0,1,3,4,5}
\pic[shift={(.3,.3)}] at ({0+(3-.6)/5*\x},{0+(2-.6)/3*\y}) {pp};

\foreach \y in {2}
\foreach \x in {0,1,2,3,5}
\pic[shift={(.3,.3)}] at ({0+(3-.6)/5*\x},{0+(2-.6)/3*\y}) {pp};

\foreach \y in {3}
\foreach \x in {1,...,5}
\pic[shift={(.3,.3)}] at ({0+(3-.6)/5*\x},{0+(2-.6)/3*\y}) {pp};

%
\foreach \y in {1}
\foreach \x in {2}
\pic[shift={(.3,.3)}] at ({0+(3-.6)/5*\x},{0+(2-.6)/3*\y}) {el};

\foreach \y in {2}
\foreach \x in {4}
\pic[shift={(.3,.3)}] at ({0+(3-.6)/5*\x},{0+(2-.6)/3*\y}) {el};

\foreach \y in {3}
\foreach \x in {0}
\pic[shift={(.3,.3)}] at ({0+(3-.6)/5*\x},{0+(2-.6)/3*\y}) {el};


%%%%%%%%%%%%%%%%%%%
%%%%%%%%%%%%%%%%%%%
%%%%%%%%%%%%%%%%%%%
% n solid 
\foreach \y in {0}
\foreach \x in {0,...,3,5}
\pic[shift={(4.3,.3)}] at ({0+(3-.6)/5*\x},{0+(2-.6)/3*\y}) {el};

\foreach \y in {1}
\foreach \x in {1,2,3,4,5}
\pic[shift={(4.3,.3)}] at ({0+(3-.6)/5*\x},{0+(2-.6)/3*\y}) {el};

\foreach \y in {2}
\foreach \x in {0,...,5}
\pic[shift={(4.3,.3)}] at ({0+(3-.6)/5*\x},{0+(2-.6)/3*\y}) {el};

\foreach \y in {3}
\foreach \x in {0,1,3,4,5}
\pic[shift={(4.3,.3)}] at ({0+(3-.6)/5*\x},{0+(2-.6)/3*\y}) {el};

%
\foreach \y in {0}
\foreach \x in {4}
\pic[shift={(4.3,.3)}] at ({0+(3-.6)/5*\x},{0+(2-.6)/3*\y}) {pp};

\foreach \y in {1}
\foreach \x in {0}
\pic[shift={(4.3,.3)}] at ({0+(3-.6)/5*\x},{0+(2-.6)/3*\y}) {pp};

\foreach \y in {3}
\foreach \x in {2}
\pic[shift={(4.3,.3)}] at ({0+(3-.6)/5*\x},{0+(2-.6)/3*\y}) {pp};





%labels
\node[anchor=south west,inner xsep=0] at (0,2) {$p$-полупроводник};
\node[anchor=south east,inner xsep=0] at (2*\sr,2) {$n$-полупроводник};



    \end{tikzpicture}
\caption{$p$--$n$-Переход}
\label{fig:p132}
\end{figure}


При образовании контакта между $p$- и $n$-полупроводниками свободные электроны (синие шары), совершая тепловое движение, проникают из $n$"=области в $p$"=область, а дырки (красные окружности)~--- из $p$"=области в $n$"=область (свободные электроны и дырки стремятся <<смешаться>> друг с другом, подобно тому как стремятся смешаться две жидкости, налитые в один сосуд).

В результате этого в пограничном слое $n$"=полупроводника имеется недостаток электронов~--- этот слой оказывается заряженным \emph{положительно} (красная область). Пограничный слой же $p$"=полупроводника становится заряженным \emph{отрицательно} (синяя область) вследствие ухода из него дырок и прихода в него электронов. Заряженные пограничные слои образуют так называемый \emph{запирающий слой} (красно"=синяя область): граница $p$--$n$"=структуры находится как бы в <<конденсаторе>>, внутреннее поле которого (запирающее поле) препятствует дальнейшему переходу свободных электронов и дырок через границу. 

Пусть $p$--$n$-структура подключена к источнику тока так, что <<плюс>> источника подан на $p$-полупроводник, а <<минус>>~--- на $n$"=полупроводник (рис.\,\ref{fig:p133}).

    \begin{figure}[H]
    \centering

    \begin{circuitikz}[font=\small,thin,
    line cap=round,line join=round,
>={Stealth[inset=0pt,round,length=3mm, angle'=20]},
vec/.style={>={Stealth[inset=0pt,round,length=3.5mm, angle'=20]}}]

    \ctikzset{bipoles/americaninductor/coils=3}


\tikzset{el/.pic={
  \begin{scope}[shift={({rnd*360}:.05)}]
  \shade[ball color=NavyBlue!70] (0,0) circle (\de);
  \end{scope}}
}


\tikzset{pp/.pic={
  \begin{scope}[shift={({rnd*360}:.05)}]
  \draw[red,thick] (0,0) circle (\de);
  \end{scope}}
}


\tikzset{em/.pic={
  \begin{scope}[shift={({rnd*360}:.025)}]
  \begin{scope}[transparency group,nearly transparent]
  \begin{scope}[transparency group,semitransparent]
  \shade[ball color=NavyBlue!70] ([shift={(0:\de cm)}]0,0) circle (\de);
  \end{scope}
  \begin{scope}[transparency group]
  \shade[ball color=NavyBlue!70] ([shift={(180:{0 cm})}]0,0) circle (\de);
  \end{scope}
  \end{scope}
  \shade[ball color=NavyBlue!70] ([shift={(180:{\de cm})}]0,0) circle (\de);
  \end{scope}}
}


\tikzset{ppm/.pic={
  \begin{scope}[shift={({rnd*360}:.025)}]
  \begin{scope}[transparency group,nearly transparent]
  \begin{scope}[transparency group,semitransparent]
  \draw[red,thick] ([shift={(180:\de cm)}]0,0) circle (\de);
  \end{scope}
  \begin{scope}[transparency group]
  \draw[red,thick] ([shift={(0:{0 cm})}]0,0) circle (\de);
  \end{scope}
  \end{scope}
  \draw[red,thick] ([shift={(0:{\de cm})}]0,0) circle (\de);
  \end{scope}}
}


\def\dn{.2}
\def\de{.1}


\def\sr{3.5}
%p n
\fill[orange, nearly transparent] (0,0) rectangle++ (\sr,2);
\fill[Green, nearly transparent] (\sr,0) rectangle++ (\sr,2);


% border charges
\fill[red,path fading=fade right] (\sr,0) rectangle++ (.25,2);
\fill[NavyBlue,path fading=fade left] (\sr,0) rectangle++ (-.25,2);


% border
\draw[gray] (\sr,0) --++ (0,2);
\draw (0,0) rectangle++ (2*\sr,2);


%%%%%%%%%%%%%%%%%%%
% p solid 
\foreach \y in {0}
\foreach \x in {0,...,5}
\pic[shift={(.3,.3)}] at ({0+(3.25-.6)/5*\x},{0+(2-.6)/3*\y}) {ppm};

\foreach \y in {1}
\foreach \x in {0,1,3,4,5}
\pic[shift={(.3,.3)}] at ({0+(3.25-.6)/5*\x},{0+(2-.6)/3*\y}) {ppm};

\foreach \y in {2}
\foreach \x in {0,1,2,3,5}
\pic[shift={(.3,.3)}] at ({0+(3.25-.6)/5*\x},{0+(2-.6)/3*\y}) {ppm};

\foreach \y in {3}
\foreach \x in {1,...,5}
\pic[shift={(.3,.3)}] at ({0+(3.25-.6)/5*\x},{0+(2-.6)/3*\y}) {ppm};

%
\foreach \y in {1}
\foreach \x in {2}
\pic[shift={(.3,.3)}] at ({0+(3.25-.6)/5*\x},{0+(2-.6)/3*\y}) {em};

\foreach \y in {2}
\foreach \x in {4}
\pic[shift={(.3,.3)}] at ({0+(3.25-.6)/5*\x},{0+(2-.6)/3*\y}) {em};

\foreach \y in {3}
\foreach \x in {0}
\pic[shift={(.3,.3)}] at ({0+(3.25-.6)/5*\x},{0+(2-.6)/3*\y}) {em};


%%%%%%%%%%%%%%%%%%%
%%%%%%%%%%%%%%%%%%%
%%%%%%%%%%%%%%%%%%%
% n solid 
\foreach \y in {0}
\foreach \x in {0,...,3,5}
\pic[shift={(4.05,.3)}] at ({0+(3.25-.6)/5*\x},{0+(2-.6)/3*\y}) {em};

\foreach \y in {1}
\foreach \x in {1,2,3,4,5}
\pic[shift={(4.05,.3)}] at ({0+(3.25-.6)/5*\x},{0+(2-.6)/3*\y}) {em};

\foreach \y in {2}
\foreach \x in {0,...,5}
\pic[shift={(4.05,.3)}] at ({0+(3.25-.6)/5*\x},{0+(2-.6)/3*\y}) {em};

\foreach \y in {3}
\foreach \x in {0,1,3,4,5}
\pic[shift={(4.05,.3)}] at ({0+(3.25-.6)/5*\x},{0+(2-.6)/3*\y}) {em};

%
\foreach \y in {0}
\foreach \x in {4}
\pic[shift={(4.05,.3)}] at ({0+(3.25-.6)/5*\x},{0+(2-.6)/3*\y}) {ppm};

\foreach \y in {1}
\foreach \x in {0}
\pic[shift={(4.05,.3)}] at ({0+(3.25-.6)/5*\x},{0+(2-.6)/3*\y}) {ppm};

\foreach \y in {3}
\foreach \x in {2}
\pic[shift={(4.05,.3)}] at ({0+(3.25-.6)/5*\x},{0+(2-.6)/3*\y}) {ppm};


%labels
\node[anchor=south west,inner xsep=0] at (0,2) {$p$-полупроводник};
\node[anchor=south east,inner xsep=0] at (2*\sr,2) {$n$-полупроводник};


%E
\draw (0,1) to[*short,*-o] ++ (-1,0) node[above] {$+$};
\draw (7,1) to[*short,*-o] ++ (1,0) node[above] {$-$};

% \node at (0,.9) [circle,draw,inner sep=0pt,minimum size=1mm,label=below:$M_1$] {};


%%%%%%%%%%%%%%%%%%%%%%%%%%%%%%%%%%%%%%%%%%%%


    \end{circuitikz}
\caption{Прямое включение $p$--$n$-перехода}
\label{fig:p133}
\end{figure}

Внешнее электрическое поле, создаваемое заряженными выводами источника (<<$+$>> и <<$-$>>) внутри структуры, направлено \emph{против} запирающего поля: на рис.\,\ref{fig:p133} внешнее поле вблизи границы направлено вправо, а запирающее поле~--- влево. Переходы основных носителей заряда (свободных электронов в $n$"=полупроводнике и дырок в $p$"=полупроводнике) через границу облегчаются: внешнее поле <<помогает>> свободным электронам и дыркам преодолеть запирающий слой (они массово перемещаются через границу контакта: свободные электроны устремляются к <<плюсу>> источника, а дырки~--- к его <<минусу>>). В этом случае сопротивление $p$--$n$"=перехода мал\'о, и \emph{ток через переход оказывается большим} (источник практически замкнут \emph{накоротко}). 

Показанный на рис.\,\ref{fig:p133} способ подключения $p$--$n$"=перехода к источнику называется \emph{прямым включением} $p$--$n$"=перехода. При \emph{обратном включении} $p$--$n$"=переход подключается наоборот: $p$"=область соединяется с <<минусом>> источника, $n$"=область~--- с <<плюсом>> (тогда \emph{ток через переход пренебрежимо мал}). 





 
\end{document}
